% GENERATED FILE. Edit canonical lessons, then run npm run generate:books.
% canonical-source-hash: fnv1a64:0cae660f43784f6c
% canonical-lessons: ML-C388-paritosikam, ML-C388-anugrahikkuka, ML-C388-asamsikkuka, ML-C388-abhinandikkuka, ML-C388-sammanikkuka, ML-R388-review-381-388

\chapter{Blessings and Gifts}
\label{ch:ml-blessings-and-gifts}
\hlchaptermodality{388}

\begin{chapteropening}
\textbf{By the end of this chapter:} \emph{I can say a gratuity, to bless, to wish someone well, to congratulate and to bestow.}

Complete the last lesson of chapter 388: I can say a gratuity, to bless, to wish someone well, to congratulate and to bestow.
\end{chapteropening}

\section[\texorpdfstring{\ml{പാരിതോഷികം} (pāritōṣikaṁ)}{പാരിതോഷികം (pāritōṣikaṁ)}]{\ml{പാരിതോഷികം} (pāritōṣikaṁ) — a gratuity; a reward}

\label{lesson:ML-C388-paritosikam}



\begin{quote}
Before the new one: say the Malayalam for to put on finery, then the Malayalam for to make ready.
\end{quote}

\subsection*{You'll want to know: \ml{പാരിതോഷികം}}
\textbf{\ml{പാരിതോഷികം}} — \emph{pāritōṣikaṁ} — \textquotedblleft{}a gratuity; a reward\textquotedblright{}.

\begin{quote}
\ml{കള്ളനെ} \ml{പിടിച്ചയാൾക്ക്} \ml{പോലീസ്} \ml{ഒരു} \ml{പാരിതോഷികം} \ml{നൽകി}.
\end{quote}

\begin{quote}
\emph{kaḷḷane piṭiccayāḷkkŭ pōlīsŭ oru pāritōṣikaṁ nalki.}
\end{quote}

\begin{quote}
The police gave a reward to the man who caught the thief.
\end{quote}

A formal word: a reward for a good deed.

\begin{cousinweb}
\textbf{\ml{പാരിതോഷികം}} (\emph{pāritōṣikaṁ}) is borrowed from Sanskrit \emph{pāritoṣika}, \textquotedblleft{}gratifying; a reward, a gratuity\textquotedblright{} (Monier-Williams).
\end{cousinweb}

\begin{grammarlens}[title={What you've built}]
One more word for telling what happened.
\end{grammarlens}

\subsection*{Your turn}
\begin{itemize}
\raggedright
  \item \emph{Say it:} \emph{pāritōṣikaṁ}
  \item \emph{Say it:} \emph{pāritōṣikaṁ}, once more
  \item \emph{Say it:} \emph{kaḷḷane piṭiccayāḷkkŭ pōlīsŭ oru pāritōṣikaṁ nalki.} — The police gave a reward to the man who caught the thief.
  \item \emph{Recall:} say the Malayalam for to put on finery, then the Malayalam for to make ready, then say \emph{pāritōṣikaṁ} again
\end{itemize}

\begin{tcolorbox}[breakable,colback=teal!4,colframe=teal!35!black,title={Before you move on}]
What does \ml{പാരിതോഷികം} mean? (\textquotedblleft{}A gratuity; a reward\textquotedblright{}.) Say it once more.
\end{tcolorbox}
\section[\texorpdfstring{\ml{അനുഗ്രഹിക്കുക} (anugrahikkuka)}{അനുഗ്രഹിക്കുക (anugrahikkuka)}]{\ml{അനുഗ്രഹിക്കുക} (anugrahikkuka) — to bless}

\label{lesson:ML-C388-anugrahikkuka}



\begin{quote}
Before the new one: say the Malayalam for to make ready, then the Malayalam for a gratuity.
\end{quote}

\subsection*{You'll want to know: \ml{അനുഗ്രഹിക്കുക}}
\textbf{\ml{അനുഗ്രഹിക്കുക}} — \emph{anugrahikkuka} — \textquotedblleft{}to bless\textquotedblright{}.

\begin{quote}
\ml{മുത്തശ്ശി} \ml{വധുവിനെയും} \ml{വരനെയും} \ml{അനുഗ്രഹിച്ചു}.
\end{quote}

\begin{quote}
\emph{muttaśśi vadhuvineyuṁ varaneyuṁ anugrahiccu.}
\end{quote}

\begin{quote}
Grandmother blessed the bride and the groom.
\end{quote}

Past: \ml{അനുഗ്രഹിച്ചു}, \emph{anugrahiccu}.

The noun is \ml{അനുഗ്രഹം}, a blessing.

\begin{cousinweb}
\textbf{\ml{അനുഗ്രഹിക്കുക}} (\emph{anugrahikkuka}) is built from Sanskrit \emph{anugraha}, \textquotedblleft{}favour, kindness\textquotedblright{} (CDIAL 314), the source of \ml{അനുഗ്രഹം} (\emph{anugrahaṁ}), with the Malayalam verb ending -\ml{ിക്കുക}.
\end{cousinweb}

\begin{grammarlens}[title={What you've built}]
One more word for telling what happened.
\end{grammarlens}

\subsection*{Your turn}
\begin{itemize}
\raggedright
  \item \emph{Say it:} \emph{anugrahikkuka}
  \item \emph{Say it:} \emph{anugrahikkuka}, once more
  \item \emph{Say it:} \emph{muttaśśi vadhuvineyuṁ varaneyuṁ anugrahiccu.} — Grandmother blessed the bride and the groom.
  \item \emph{Recall:} say the Malayalam for to make ready, then the Malayalam for a gratuity, then say \emph{anugrahikkuka} again
\end{itemize}

\begin{tcolorbox}[breakable,colback=teal!4,colframe=teal!35!black,title={Before you move on}]
What does \ml{അനുഗ്രഹിക്കുക} mean? (\textquotedblleft{}To bless\textquotedblright{}.) Say it once more.
\end{tcolorbox}
\section[\texorpdfstring{\ml{ആശംസിക്കുക} (āśaṁsikkuka)}{ആശംസിക്കുക (āśaṁsikkuka)}]{\ml{ആശംസിക്കുക} (āśaṁsikkuka) — to wish someone well}

\label{lesson:ML-C388-asamsikkuka}



\begin{quote}
Before the new one: say the Malayalam for a gratuity, then the Malayalam for to bless.
\end{quote}

\subsection*{You'll want to know: \ml{ആശംസിക്കുക}}
\textbf{\ml{ആശംസിക്കുക}} — \emph{āśaṁsikkuka} — \textquotedblleft{}to wish someone well\textquotedblright{}.

\begin{quote}
\ml{എല്ലാവരും} \ml{ദമ്പതികൾക്ക്} \ml{നല്ല} \ml{ഭാവി} \ml{ആശംസിച്ചു}.
\end{quote}

\begin{quote}
\emph{ellāvaruṁ dampatikaḷkkŭ nalla bhāvi āśaṁsiccu.}
\end{quote}

\begin{quote}
Everyone wished the couple a good future.
\end{quote}

Past: \ml{ആശംസിച്ചു}, \emph{āśaṁsiccu}.

The person takes -\ml{ക്ക്} or -\ml{ന്}. The noun is \ml{ആശംസകൾ}, good wishes.

\begin{cousinweb}
\textbf{\ml{ആശംസിക്കുക}} (\emph{āśaṁsikkuka}) is built from the Sanskrit verb \emph{ā-śaṁs}, \textquotedblleft{}to hope for, to wish\textquotedblright{} (Monier-Williams), whose noun gave \ml{ആശംസകൾ} (\emph{āśaṁsakaḷ}), with the Malayalam verb ending -\ml{ിക്കുക}.
\end{cousinweb}

\begin{grammarlens}[title={What you've built}]
One more word for telling what happened.
\end{grammarlens}

\subsection*{Your turn}
\begin{itemize}
\raggedright
  \item \emph{Say it:} \emph{āśaṁsikkuka}
  \item \emph{Say it:} \emph{āśaṁsikkuka}, once more
  \item \emph{Say it:} \emph{ellāvaruṁ dampatikaḷkkŭ nalla bhāvi āśaṁsiccu.} — Everyone wished the couple a good future.
  \item \emph{Recall:} say the Malayalam for a gratuity, then the Malayalam for to bless, then say \emph{āśaṁsikkuka} again
\end{itemize}

\begin{tcolorbox}[breakable,colback=teal!4,colframe=teal!35!black,title={Before you move on}]
What does \ml{ആശംസിക്കുക} mean? (\textquotedblleft{}To wish someone well\textquotedblright{}.) Say it once more.
\end{tcolorbox}
\section[\texorpdfstring{\ml{അഭിനന്ദിക്കുക} (abhinandikkuka)}{അഭിനന്ദിക്കുക (abhinandikkuka)}]{\ml{അഭിനന്ദിക്കുക} (abhinandikkuka) — to congratulate}

\label{lesson:ML-C388-abhinandikkuka}



\begin{quote}
Before the new one: say the Malayalam for to bless, then the Malayalam for to wish someone well.
\end{quote}

\subsection*{You'll want to know: \ml{അഭിനന്ദിക്കുക}}
\textbf{\ml{അഭിനന്ദിക്കുക}} — \emph{abhinandikkuka} — \textquotedblleft{}to congratulate\textquotedblright{}.

\begin{quote}
\ml{അധ്യാപകൻ} \ml{എന്നെ} \ml{നല്ല} \ml{മാർക്കിന്} \ml{അഭിനന്ദിച്ചു}.
\end{quote}

\begin{quote}
\emph{adhyāpakan enne nalla mārkkinŭ abhinandiccu.}
\end{quote}

\begin{quote}
The teacher congratulated me on my good marks.
\end{quote}

Past: \ml{അഭിനന്ദിച്ചു}, \emph{abhinandiccu}.

The person takes -\ml{എ}, the reason -\ml{ക്ക്} or -\ml{ന്}. \ml{അഭിനന്ദനങ്ങൾ}! is Congratulations!

\begin{cousinweb}
\textbf{\ml{അഭിനന്ദിക്കുക}} (\emph{abhinandikkuka}) is built from the Sanskrit verb \emph{abhi-nand}, \textquotedblleft{}to rejoice at, to greet, to praise\textquotedblright{} (Monier-Williams), with the Malayalam verb ending -\ml{ിക്കുക}.
\end{cousinweb}

\begin{grammarlens}[title={What you've built}]
One more word for telling what happened.
\end{grammarlens}

\subsection*{Your turn}
\begin{itemize}
\raggedright
  \item \emph{Say it:} \emph{abhinandikkuka}
  \item \emph{Say it:} \emph{abhinandikkuka}, once more
  \item \emph{Say it:} \emph{adhyāpakan enne nalla mārkkinŭ abhinandiccu.} — The teacher congratulated me on my good marks.
  \item \emph{Recall:} say the Malayalam for to bless, then the Malayalam for to wish someone well, then say \emph{abhinandikkuka} again
\end{itemize}

\begin{tcolorbox}[breakable,colback=teal!4,colframe=teal!35!black,title={Before you move on}]
What does \ml{അഭിനന്ദിക്കുക} mean? (\textquotedblleft{}To congratulate\textquotedblright{}.) Say it once more.
\end{tcolorbox}
\section[\texorpdfstring{\ml{സമ്മാനിക്കുക} (sammānikkuka)}{സമ്മാനിക്കുക (sammānikkuka)}]{\ml{സമ്മാനിക്കുക} (sammānikkuka) — to bestow; to give as a gift}

\label{lesson:ML-C388-sammanikkuka}



\begin{quote}
Before the new one: say the Malayalam for to wish someone well, then the Malayalam for to congratulate.
\end{quote}

\subsection*{You'll want to know: \ml{സമ്മാനിക്കുക}}
\textbf{\ml{സമ്മാനിക്കുക}} — \emph{sammānikkuka} — \textquotedblleft{}to bestow; to give as a gift\textquotedblright{}.

\begin{quote}
\ml{അച്ഛൻ} \ml{അമ്മയ്ക്ക്} \ml{ഒരു} \ml{മോതിരം} \ml{സമ്മാനിച്ചു}.
\end{quote}

\begin{quote}
\emph{acchan ammaykkŭ oru mōtiraṁ sammāniccu.}
\end{quote}

\begin{quote}
Father gave mother a ring as a gift.
\end{quote}

Past: \ml{സമ്മാനിച്ചു}, \emph{sammāniccu}.

From \ml{സമ്മാനം}, a gift, a prize.

\begin{cousinweb}
\textbf{\ml{സമ്മാനിക്കുക}} (\emph{sammānikkuka}) is built on \ml{സമ്മാനം} (\emph{sammānaṁ}), \textquotedblleft{}a gift, a prize\textquotedblright{}, from Sanskrit \emph{sammāna}, \textquotedblleft{}honour, respect\textquotedblright{} (Monier-Williams), with the Malayalam verb ending -\ml{ിക്കുക}.
\end{cousinweb}

\begin{grammarlens}[title={What you've built}]
One more word for telling what happened.
\end{grammarlens}

\subsection*{Your turn}
\begin{itemize}
\raggedright
  \item \emph{Say it:} \emph{sammānikkuka}
  \item \emph{Say it:} \emph{sammānikkuka}, once more
  \item \emph{Say it:} \emph{acchan ammaykkŭ oru mōtiraṁ sammāniccu.} — Father gave mother a ring as a gift.
  \item \emph{Recall:} say the Malayalam for to wish someone well, then the Malayalam for to congratulate, then say \emph{sammānikkuka} again
\end{itemize}

\begin{tcolorbox}[breakable,colback=teal!4,colframe=teal!35!black,title={Before you move on}]
What does \ml{സമ്മാനിക്കുക} mean? (\textquotedblleft{}To bestow; to give as a gift\textquotedblright{}.) Say it once more.
\end{tcolorbox}
\section[(dialogue)]{Words from chapters 381–388 — a review}

\label{lesson:ML-R388-review-381-388}



\begin{quote}
Say the Malayalam for darkness and a sound.
\end{quote}

\subsection*{Your turn}
Say these aloud:

\begin{itemize}
\raggedright
  \item darkness, a sound, silence, to whisper, to disappear
  \item a riot, a weapon, to surround, to drive away, to calm down
  \item a secret, to appear, to vanish, to find out, to retrieve
  \item birth, death, an obituary, an autobiography, a biography
  \item a marriage, an engagement, a wedding pendant, an invitation card, to marry
  \item a host, hospitality, a reception, to invite, to receive
  \item a banquet, a house guest, a ceremony, to put on finery, to make ready
  \item a gratuity, to bless, to wish someone well, to congratulate, to bestow
\end{itemize}

\begin{tcolorbox}[breakable,colback=teal!4,colframe=teal!35!black,title={Before you move on}]
Darkness? (\textbf{\ml{ഇരുട്ട്}}, \emph{iruṭṭŭ}.) A sound? (\textbf{\ml{ശബ്ദം}}, \emph{śabdaṁ}.) Silence? (\textbf{\ml{നിശ്ശബ്ദത}}, \emph{niśśabdata}.) To whisper? (\textbf{\ml{മന്ത്രിക്കുക}}, \emph{mantrikkuka}.) To disappear? (\textbf{\ml{അപ്രത്യക്ഷമാകുക}}, \emph{apratyakṣamākuka}.) A riot? (\textbf{\ml{കലാപം}}, \emph{kalāpaṁ}.) A weapon? (\textbf{\ml{ആയുധം}}, \emph{āyudhaṁ}.) To surround? (\textbf{\ml{വളയുക}}, \emph{vaḷayuka}.) To drive away? (\textbf{\ml{തുരത്തുക}}, \emph{turattuka}.) To calm down? (\textbf{\ml{ശാന്തമാകുക}}, \emph{śāntamākuka}.) A secret? (\textbf{\ml{രഹസ്യം}}, \emph{rahasyaṁ}.) To appear? (\textbf{\ml{പ്രത്യക്ഷപ്പെടുക}}, \emph{pratyakṣappeṭuka}.) To vanish? (\textbf{\ml{മറയുക}}, \emph{maṟayuka}.) To find out? (\textbf{\ml{കണ്ടുപിടിക്കുക}}, \emph{kaṇṭupiṭikkuka}.) To retrieve? (\textbf{\ml{കണ്ടെടുക്കുക}}, \emph{kaṇṭeṭukkuka}.) Birth? (\textbf{\ml{ജനനം}}, \emph{jananaṁ}.) Death? (\textbf{\ml{മരണം}}, \emph{maraṇaṁ}.) An obituary? (\textbf{\ml{ചരമം}}, \emph{caramaṁ}.) An autobiography? (\textbf{\ml{ആത്മകഥ}}, \emph{ātmakatha}.) A biography? (\textbf{\ml{ജീവചരിത്രം}}, \emph{jīvacaritraṁ}.) A marriage? (\textbf{\ml{വിവാഹം}}, \emph{vivāhaṁ}.) An engagement? (\textbf{\ml{വിവാഹനിശ്ചയം}}, \emph{vivāhaniścayaṁ}.) A wedding pendant? (\textbf{\ml{താലി}}, \emph{tāli}.) An invitation card? (\textbf{\ml{ക്ഷണക്കത്ത്}}, \emph{kṣaṇakkattŭ}.) To marry? (\textbf{\ml{വിവാഹം} \ml{കഴിക്കുക}}, \emph{vivāhaṁ kaḻikkuka}.) A host? (\textbf{\ml{ആതിഥേയൻ}}, \emph{ātithēyan}.) Hospitality? (\textbf{\ml{ആതിഥ്യം}}, \emph{ātithyaṁ}.) A reception? (\textbf{\ml{സ്വീകരണം}}, \emph{svīkaraṇaṁ}.) To invite? (\textbf{\ml{ക്ഷണിക്കുക}}, \emph{kṣaṇikkuka}.) To receive? (\textbf{\ml{സ്വീകരിക്കുക}}, \emph{svīkarikkuka}.) A banquet? (\textbf{\ml{സൽക്കാരം}}, \emph{salkkāraṁ}.) A house guest? (\textbf{\ml{വിരുന്നുകാരൻ}}, \emph{virunnukāran}.) A ceremony? (\textbf{\ml{ചടങ്ങ്}}, \emph{caṭaṅṅŭ}.) To put on finery? (\textbf{\ml{അണിയുക}}, \emph{aṇiyuka}.) To make ready? (\textbf{\ml{ഒരുക്കുക}}, \emph{orukkuka}.) A gratuity? (\textbf{\ml{പാരിതോഷികം}}, \emph{pāritōṣikaṁ}.) To bless? (\textbf{\ml{അനുഗ്രഹിക്കുക}}, \emph{anugrahikkuka}.) To wish someone well? (\textbf{\ml{ആശംസിക്കുക}}, \emph{āśaṁsikkuka}.) To congratulate? (\textbf{\ml{അഭിനന്ദിക്കുക}}, \emph{abhinandikkuka}.) To bestow? (\textbf{\ml{സമ്മാനിക്കുക}}, \emph{sammānikkuka}.)
\end{tcolorbox}
